\documentclass[11pt,a4paper]{article}

\usepackage{amsmath,amssymb,amsthm}
\usepackage{bm}
\usepackage{mathtools}
\usepackage{physics}
\usepackage[colorlinks=true,linkcolor=blue,urlcolor=blue,citecolor=blue]{hyperref}

\DeclareMathOperator{\E}{\mathbb{E}}
\DeclareMathOperator{\svd}{SVD}
\DeclareMathOperator{\tr}{tr}

\title{LoRC: Low-rank Quantization Correction}
\date{}

\begin{document}
\maketitle

\section{Introduction}

We consider the problem of compressing a pretrained language model
$f_\theta$ while preserving its performance on a target domain
$\mathcal{D}_{\mathrm{target}}$, potentially at the expense of
out-of-domain generality. Our approach---\textbf{Lo}w-rank
\textbf{Q}uantization \textbf{C}orrection---identifies low-dimensional
subspaces in which activation statistics differ between domains and
uses these subspaces to construct low-rank corrections to quantized
weights.

The method is motivated by the observation that quantization error is
not equally important in all activation directions. If a domain
frequently occupies a particular subspace, errors in the corresponding
weight directions can have a disproportionately large effect on the
domain's output.

\subsection{Covariance estimation}

For a linear layer with weight
\[
W \in \mathbb{R}^{d_{\mathrm{out}} \times d_{\mathrm{in}}},
\]
we use column-vector activations
$x \in \mathbb{R}^{d_{\mathrm{in}}}$ and
$y \in \mathbb{R}^{d_{\mathrm{out}}}$, with
\[
y = Wx.
\]

For each domain $\mathcal{D}$, we estimate the covariance matrices
\begin{align}
C_{\mathrm{pre}}^{\mathcal{D}}
&=
\E[xx^{\mathsf T}]
-
\E[x]\E[x]^{\mathsf T}
\in
\mathbb{R}^{d_{\mathrm{in}}\times d_{\mathrm{in}}},
\\
C_{\mathrm{post}}^{\mathcal{D}}
&=
\E[yy^{\mathsf T}]
-
\E[y]\E[y]^{\mathsf T}
\in
\mathbb{R}^{d_{\mathrm{out}}\times d_{\mathrm{out}}}.
\end{align}

Given samples $\{x_i\}_{i=1}^N$ from a domain, the corresponding
population-normalized estimator is
\begin{align}
\widehat{\mu}
&=
\frac{1}{N}\sum_{i=1}^N x_i,
\\
\widehat{C}
&=
\frac{1}{N}\sum_{i=1}^N x_i x_i^{\mathsf T}
-
\widehat{\mu}\widehat{\mu}^{\mathsf T}.
\end{align}

The normalization is performed independently for each domain so that
differences in the number of collected samples do not affect the
estimated covariance.

\subsection{Domain subspace extraction}

Given two domains $\mathcal{D}_1$ and $\mathcal{D}_2$, we form the
difference between their covariance matrices:
\begin{align}
\Delta C
&=
C^{\mathcal{D}_1}
-
C^{\mathcal{D}_2}.
\end{align}

In exact arithmetic, both covariance matrices are symmetric, and
therefore $\Delta C$ is symmetric. For numerical estimates, we
explicitly symmetrize:
\begin{align}
\Delta C
&\leftarrow
\frac{1}{2}
\left(
\Delta C+\Delta C^{\mathsf T}
\right).
\end{align}

More generally, weighted covariances can be used:
\begin{align}
\Delta C
&=
\alpha C^{\mathcal{D}_1}
-
\beta C^{\mathcal{D}_2},
\qquad
\alpha,\beta\geq 0.
\end{align}
In this case the eigenvalues measure differences in
\emph{weighted} activation variance. Unless there is a specific
normalization objective, we use $\alpha=\beta=1$.

We compute the eigendecomposition
\begin{align}
\Delta C
&=
Q\Lambda Q^{\mathsf T},
\end{align}
where $Q$ is orthogonal and $\Lambda$ is diagonal.

For any unit vector $v$,
\begin{align}
v^{\mathsf T}\Delta C v
&=
v^{\mathsf T}C^{\mathcal{D}_1}v
-
v^{\mathsf T}C^{\mathcal{D}_2}v.
\end{align}

Thus eigenvectors associated with the largest positive eigenvalues
identify directions with substantially greater activation variance in
$\mathcal{D}_1$, while eigenvectors associated with the most negative
eigenvalues identify directions with substantially greater activation
variance in $\mathcal{D}_2$.

We retain the $K$ largest positive eigen-directions and the $K$
most-negative eigen-directions, subject to their availability, forming
column-orthonormal matrices
\[
V_1,V_2\in\mathbb{R}^{d\times K}.
\]

The resulting subspaces describe differences in activation covariance;
they do not by themselves imply that quantization error is concentrated
in those directions.

\subsection{Connection to quantization error}

Let $W_{\mathrm{base}}$ denote the dequantized NF4 weight and define
the quantization error
\begin{align}
E
&=
W-W_{\mathrm{base}},
\\
W_{\mathrm{base}}
&=
Q^{-1}(Q(W)).
\end{align}

For a centered input activation $x$ with covariance $C$, the output
error caused by $E$ is
\[
\delta y = Ex.
\]
Its expected squared norm is
\begin{align}
\E\left[\|\delta y\|_2^2\right]
&=
\E\left[x^{\mathsf T}E^{\mathsf T}Ex\right]
\\
&=
\tr\left(ECE^{\mathsf T}\right).
\end{align}

Consequently, the difference in expected quantization-error energy
between two domains is
\begin{align}
\Delta L_E
&=
\tr\left(
E
(C^{\mathcal{D}_1}-C^{\mathcal{D}_2})
E^{\mathsf T}
\right)
\\
&=
\tr\left(E\Delta C E^{\mathsf T}\right).
\end{align}

This provides a direct connection between domain-dependent activation
statistics and the effect of quantization. In particular, activation
subspaces are useful for correction when the corresponding directions
also contain significant quantization error.

\subsection{Correction factors}

Consider first an input-space domain subspace
\[
V\in\mathbb{R}^{d_{\mathrm{in}}\times K},
\qquad
V^{\mathsf T}V=I.
\]

The component of the quantization error acting on this subspace is
\begin{align}
E_{\mathrm{proj}}
&=
E(VV^{\mathsf T}).
\end{align}

Rather than constructing the full
$d_{\mathrm{in}}\times d_{\mathrm{in}}$ projection matrix, define
\begin{align}
A
&=
EV
\in
\mathbb{R}^{d_{\mathrm{out}}\times K}.
\end{align}

Then
\begin{align}
E_{\mathrm{proj}}
&=
AV^{\mathsf T}.
\end{align}

Compute the compact SVD
\begin{align}
A
&=
U_A\Sigma_A R_A^{\mathsf T}.
\end{align}

Keeping the first $k$ singular components, where
\[
k=\min(K,d_{\mathrm{out}},d_{\mathrm{in}}),
\]
gives
\begin{align}
A
&\approx
U_{A,k}\Sigma_{A,k}R_{A,k}^{\mathsf T}.
\end{align}

Therefore
\begin{align}
E_{\mathrm{proj}}
&\approx
U_{A,k}\Sigma_{A,k}
R_{A,k}^{\mathsf T}V^{\mathsf T}
\\
&=
U_{A,k}\Sigma_{A,k}
(VR_{A,k})^{\mathsf T}.
\end{align}

Define
\begin{align}
V_{\mathrm{act}}
&=
(VR_{A,k})\Sigma_{A,k}^{1/2}
\in
\mathbb{R}^{d_{\mathrm{in}}\times k},
\\
U_{\mathrm{write}}
&=
U_{A,k}\Sigma_{A,k}^{1/2}
\in
\mathbb{R}^{d_{\mathrm{out}}\times k}.
\end{align}

The resulting low-rank approximation satisfies
\begin{align}
E_{\mathrm{proj}}
&\approx
U_{\mathrm{write}}V_{\mathrm{act}}^{\mathsf T}.
\end{align}

The factors are stored in bfloat16.

For an output-space domain subspace
\[
V\in\mathbb{R}^{d_{\mathrm{out}}\times K},
\]
the corresponding projection is
\begin{align}
E_{\mathrm{proj}}
&=
(VV^{\mathsf T})E.
\end{align}

The analogous low-rank factorization can be obtained by applying the
same procedure to $E^{\mathsf T}$ and transposing the resulting
factors.

\subsection{Hybrid forward pass}

For a row-vector representation of the input activation,
\[
x\in\mathbb{R}^{1\times d_{\mathrm{in}}},
\]
the base forward pass is
\begin{align}
z_{\mathrm{base}}
&=
xW_{\mathrm{base}}^{\mathsf T}.
\end{align}

Using the low-rank approximation
\[
E_{\mathrm{proj}}
\approx
U_{\mathrm{write}}V_{\mathrm{act}}^{\mathsf T},
\]
the corresponding output correction is
\begin{align}
z_{\mathrm{corr}}
&=
xE_{\mathrm{proj}}^{\mathsf T}
\\
&\approx
xV_{\mathrm{act}}U_{\mathrm{write}}^{\mathsf T}
\\
&=
(xV_{\mathrm{act}})
U_{\mathrm{write}}^{\mathsf T}.
\end{align}

Thus the hybrid forward pass is
\begin{align}
\boxed{
z
=
xW_{\mathrm{base}}^{\mathsf T}
+
(xV_{\mathrm{act}})
U_{\mathrm{write}}^{\mathsf T}.
}
\end{align}

The first term uses the quantized NF4 weight, while the second term
uses the small bfloat16 low-rank correction.

\subsection{Spectral rank selection}

The singular values of
\[
A=EV
\]
measure the strength of the best orthogonal low-rank components of the
projected quantization error. They should not be interpreted as
contributions of the individual original domain eigenvectors, because
the right singular vectors $R_A$ can rotate and mix those directions.

A simple rank-selection rule is therefore to retain singular values
above a threshold:
\begin{align}
\mathrm{keep}_i
&=
\mathbf{1}(\sigma_i>\tau).
\end{align}

If a fixed fraction $\mathrm{keep\_pct}$ of the singular components is
desired, the threshold can be selected as
\begin{align}
\tau
&=
\operatorname{quantile}
\left(
\{\sigma_i\},
1-\mathrm{keep\_pct}
\right).
\end{align}

Alternatively, for a fixed rank $k$, the $k$ largest singular values
are retained directly.

This procedure provides a deterministic spectral rank-selection
heuristic and does not require gradient-based mask training.

\subsection{Disjunction score}

To evaluate whether a domain subspace is genuinely domain-specific,
we can measure the effect of removing the corresponding subspace from
the original weight matrix.

For an input-space subspace $V$,
\begin{align}
W'
&=
W-W(VV^{\mathsf T}).
\end{align}

For an output-space subspace $V$,
\begin{align}
W'
&=
W-(VV^{\mathsf T})W.
\end{align}

Let $\Delta\mathrm{PPL}_{\mathrm{own}}$ denote the perplexity change
on the domain from which the subspace was extracted, and let
$\Delta\mathrm{PPL}_{\mathrm{other}}$ denote the corresponding change
on another domain.

A clipped ratio can be defined as
\begin{align}
D
&=
\frac{
\max(\Delta\mathrm{PPL}_{\mathrm{own}},\varepsilon)
}{
\max(\Delta\mathrm{PPL}_{\mathrm{other}},\varepsilon)
},
\qquad
\varepsilon>0.
\end{align}

Large $D$ indicates that the subspace has a substantially larger
effect on its source domain than on the comparison domain.

This diagnostic measures domain dependence of the underlying weight
directions. It should be distinguished from a direct measurement of
the effectiveness of the quantization correction itself. For the
latter, one should also measure the change obtained by removing
\[
U_{\mathrm{write}}V_{\mathrm{act}}^{\mathsf T}
\]
from the quantized model.

\subsection{Streaming covariance products}

Explicitly storing a dense covariance matrix requires
$\mathcal{O}(d^2)$ memory. For large transformer layers, it is preferable
to estimate covariance--vector products directly.

For a domain with mean $\mu=\E[x]$,
\begin{align}
Cv
&=
\E[x(x^{\mathsf T}v)]
-
\mu(\mu^{\mathsf T}v).
\end{align}

This permits iterative eigensolvers to operate without explicitly
constructing $C$.

A low-rank subspace estimate can alternatively be maintained using
online PCA methods. For example, a simple Oja update is
\begin{align}
V
&\leftarrow
V+\eta\,x(x^{\mathsf T}V),
\end{align}
followed by orthonormalization,
\begin{align}
V
&\leftarrow
\operatorname{QR}(V).
\end{align}

This reduces persistent storage from
$\mathcal{O}(d^2)$ to $\mathcal{O}(dK)$, although the computational cost
of processing the activation stream remains dependent on $dK$ per
sample.

For the two-domain difference, separate covariance products can be
estimated and combined as
\begin{align}
\Delta C\,v
&=
C^{\mathcal{D}_1}v
-
C^{\mathcal{D}_2}v.
\end{align}

\subsection{Randomized and Krylov eigendecomposition}

Because $\Delta C$ is symmetric but generally indefinite, an eigensolver
must be capable of recovering both large positive and large negative
eigenvalues.

A Lanczos method is a natural choice because it operates using only
matrix--vector products with $\Delta C$.

A randomized range finder can also be used. Draw
\begin{align}
\Omega
&\sim
\mathcal{N}(0,1)^{d\times(K+p)}
\end{align}
and compute
\begin{align}
Y
&=
\Delta C\,\Omega.
\end{align}

For improved capture of both ends of the spectrum, optional power
iterations can be applied:
\begin{align}
Y
&=
(\Delta C)^{2q+1}\Omega.
\end{align}

Since $\Delta C$ is symmetric, the magnitude of an eigenvalue
$\lambda_i$ is amplified as
\[
|\lambda_i|^{2q+1},
\]
allowing both large positive and large negative eigenvalues to be
captured.

After orthonormalization,
\begin{align}
Q
&=
\operatorname{QR}(Y),
\end{align}
form the Rayleigh--Ritz matrix
\begin{align}
T
&=
Q^{\mathsf T}\Delta C Q.
\end{align}

Then compute
\begin{align}
T
&=
\widetilde U
\widetilde\Lambda
\widetilde U^{\mathsf T},
\end{align}
and approximate the corresponding eigenvectors by
\begin{align}
V_{\mathrm{ext}}
&=
Q\widetilde U.
\end{align}

For dense covariance operators, each matrix--block multiplication
costs approximately
\[
\mathcal{O}(d^2(K+p)).
\]
Thus the total cost depends on the number of Krylov or power iterations,
but avoids the $\mathcal{O}(d^3)$ cost of a full dense
eigendecomposition.

\subsection{Complexity of the low-rank correction}

For an input-space subspace $V\in\mathbb{R}^{d_{\mathrm{in}}\times K}$,
forming
\[
A=EV
\]
costs
\[
\mathcal{O}(d_{\mathrm{out}}d_{\mathrm{in}}K)
\]
for a dense error matrix.

The subsequent SVD of
\[
A\in\mathbb{R}^{d_{\mathrm{out}}\times K}
\]
is substantially cheaper than an SVD of the full
$d_{\mathrm{out}}\times d_{\mathrm{in}}$ matrix. For
$K\ll d_{\mathrm{in}},d_{\mathrm{out}}$, its cost is approximately
\[
\mathcal{O}(d_{\mathrm{out}}K^2)
\]
using a standard dense thin SVD.

The projection matrix $VV^{\mathsf T}$ therefore never needs to be
constructed.

If the quantized weight is processed blockwise, the full-precision
error matrix $E$ also need not be materialized, allowing the correction
to be computed directly from quantized-weight blocks.

\subsection{Multi-domain extension}

For $D>2$ domains, pairwise covariance differences can be generalized
using Fisher's linear discriminant framework.

Let
\[
\mu_d
=
\E[x\mid\mathcal{D}_d],
\qquad
\mu
=
\sum_{d=1}^D\pi_d\mu_d,
\]
where $\pi_d$ is the empirical probability of domain $d$.

The between-domain covariance is
\begin{align}
C_{\mathrm{between}}
&=
\sum_{d=1}^D
\pi_d
(\mu_d-\mu)(\mu_d-\mu)^{\mathsf T}.
\end{align}

The pooled within-domain covariance is
\begin{align}
C_{\mathrm{within}}
&=
\sum_{d=1}^D
\pi_d C_d.
\end{align}

A generalized Fisher subspace can be obtained from
\begin{align}
C_{\mathrm{between}}v
&=
\lambda
C_{\mathrm{within}}v.
\end{align}

Equivalently, one can solve
\begin{align}
\max_{W}
\quad
&\tr(W^{\mathsf T}C_{\mathrm{between}}W)
\\
\text{subject to}
\quad
&W^{\mathsf T}C_{\mathrm{within}}W=I.
\end{align}

In high-dimensional activation spaces, $C_{\mathrm{within}}$ may be
singular or poorly conditioned. A regularized version is therefore
preferable:
\begin{align}
C_{\mathrm{within}}^{(\lambda)}
&=
C_{\mathrm{within}}+\lambda I,
\qquad
\lambda>0.
\end{align}

The resulting generalized eigenvectors provide a single subspace that
captures domain-discriminative mean variation. If covariance
differences rather than mean differences are the primary signal of
interest, a separate covariance-discriminant construction is required;
standard Fisher LDA alone does not capture all such covariance
differences.

\subsection{Cross-layer subspace coupling}

Per-module subspaces can be estimated independently. To reduce
overfitting and parameter count, modules with compatible activation
dimensions can share a global basis
\[
V_{\mathrm{gl}}\in\mathbb{R}^{d\times K}.
\]

A module-specific subspace can then be represented as
\begin{align}
V_\ell
&=
V_{\mathrm{gl}}M_\ell.
\end{align}

If $V_{\mathrm{gl}}$ is column-orthonormal and $V_\ell$ is required to
remain column-orthonormal, then $M_\ell$ should satisfy
\begin{align}
M_\ell^{\mathsf T}M_\ell
&=
I.
\end{align}

This construction applies directly only to modules sharing the same
activation dimension $d$. For modules with different dimensions,
separate global bases or learned mappings into a common latent space
are required.

\subsection{Practical feasibility}

The full covariance implementation requires
$\mathcal{O}(d^2)$ storage and can require $\mathcal{O}(Nd^2)$ work for
$N$ activation samples. This is practical for small models and
selected layers, but becomes expensive for large language models.

The scalable implementation should therefore use covariance--vector
products or streaming low-rank PCA, requiring approximately
$\mathcal{O}(dK)$ persistent subspace storage.

The low-rank correction itself is considerably cheaper. Once the
domain subspace $V$ has been obtained, the matrix
\[
A=EV
\]
has only $K$ columns, and its compact SVD yields the correction factors
without constructing the full projected error.

A useful diagnostic before deploying the correction is the fraction of
quantization-error energy captured by the selected subspace:
\begin{align}
r
&=
\frac{
\|EVV^{\mathsf T}\|_F^2
}{
\|E\|_F^2
}.
\end{align}

For comparison against a random $K$-dimensional subspace, one can
measure the corresponding ratio
\[
r_{\mathrm{random}}.
\]

If
\[
r\gg r_{\mathrm{random}},
\]
the domain-derived subspace is capturing unusually large amounts of
quantization error. This provides a stronger sanity check than
activation-domain separation alone.

Finally, the correction should be evaluated against the following
baselines:
\begin{enumerate}
  \item the original full-precision model;
  \item the NF4-quantized model without correction;
  \item a random low-rank correction;
  \item a domain-derived low-rank correction;
  \item a directly optimized low-rank correction.
\end{enumerate}

The central empirical criterion is whether a small number of bfloat16
parameters can recover a substantial fraction of the target-domain
performance lost under NF4 while providing substantially less benefit
on the comparison domain.

\end{document}